\documentclass[a4paper]{arrowhead}

\usepackage[yyyymmdd]{datetime}
\usepackage{etoolbox}
\usepackage[T1]{fontenc}
\usepackage[utf8]{inputenc}
\usepackage{multirow}
\usepackage{longtable}
\usepackage{amssymb}
\usepackage{float}
\usepackage{tikz}
\usetikzlibrary{positioning,arrows.meta,fit,calc,shapes.misc,shapes.geometric}

\renewcommand{\dateseparator}{-}
\setlength{\emergencystretch}{3em}
\lstset{basicstyle=\ttfamily\footnotesize,breaklines=true,columns=fullflexible,frame=single,rulecolor=\color{gray!50},backgroundcolor=\color{gray!6}}
\tikzset{
  ahsys/.style={draw=ArrowheadBlue,very thick,fill=ArrowheadBlue!6,rounded corners=2pt,minimum width=3.2cm,minimum height=1.1cm,align=center,font=\small\bfseries},
  ahext/.style={draw=ArrowheadGrey,thick,dashed,fill=gray!5,rounded corners=2pt,minimum width=3.2cm,minimum height=1.0cm,align=center,font=\small},
  ahsrv/.style={draw=ArrowheadPurple,thick,fill=ArrowheadPurple!8,rounded corners=6pt,minimum height=0.55cm,inner xsep=4pt,align=center,font=\scriptsize\ttfamily},
  ahcsm/.style={draw=ArrowheadTeal,thick,fill=ArrowheadTeal!8,rounded corners=6pt,minimum height=0.55cm,inner xsep=4pt,align=center,font=\scriptsize\ttfamily},
  ahdb/.style={draw=ArrowheadOrange,thick,fill=ArrowheadOrange!8,cylinder,shape border rotate=90,aspect=0.25,minimum height=0.9cm,align=center,font=\scriptsize},
  ahflow/.style={-{Stealth[length=2.2mm]},thick,ArrowheadGrey},
  ahlabel/.style={font=\scriptsize,fill=white,inner sep=1pt,align=center},
}

%% Special references
\newcommand{\fref}[1]{{\textcolor{ArrowheadBlue}{\hyperref[sec:functions:#1]{#1}}}}
\newcommand{\mref}[1]{{\textcolor{ArrowheadPurple}{\hyperref[sec:model:#1]{#1}}}}
\newcommand{\sref}[1]{{\textcolor{ArrowheadPurple}{#1}}}
\newrobustcmd\msubsection[2]{
  \addtocounter{subsection}{1}
  \addcontentsline{toc}{subsection}{\protect\numberline{\thesubsection}#1 \textcolor{ArrowheadPurple}{#2}}
  \subsection*{\thesubsection\quad#1 \textcolor{ArrowheadPurple}{#2}}
  \label{sec:model:#2}
}

\begin{document}

%% Arrowhead Document Properties
\ArrowheadTitle{WorkflowManager}

\ArrowheadType{System Description}
\ArrowheadTypeShort{SysD}
\ArrowheadVersion{5.4.0}
\ArrowheadDate{2026-11-23}
\ArrowheadAuthor{Documentation maintenance}
\ArrowheadStatus{REVIEW DRAFT}
\ArrowheadContact{Repository maintainers}
\ArrowheadFooter{\href{www.arrowhead.eu}{www.arrowhead.eu}}
\ArrowheadSetup
%%

%% Front Page
\begin{center}
  \vspace*{1cm}
  \huge{\arrowtitle}

  \vspace*{0.2cm}
  \LARGE{\arrowtype}
  \vspace*{1cm}

  
  \vspace*{\fill}

  

  \vspace*{1cm}
  \vspace*{\fill}

  % Front Page Abstract
  \begin{abstract}
    This document describes the WorkflowManager, an Eclipse Arrowhead support system that keeps the books of workflows: it holds recipes and sessions, validates recipes, advances sessions with the transition rules of the formal model, dispatches steps to WorkflowExecutors and reports to the product. Documentation edition 5.4.0; system version 5.4.0 (framework), service versions 1.0.0.
  \end{abstract}

  \vspace*{1cm}
\end{center}

\newpage
%%

%% Table of Contents
\tableofcontents
\newpage
%%

\section{Overview}
\label{sec:overview}
The WorkflowManager (WfM) is the book-keeper of the Eclipse Arrowhead workflow support systems. It
stores workflow \emph{recipes} (which services are consumed, in which order and under which conditions), runs
\emph{sessions} of them for products, control periods or batches, and keeps the session state of the formal execution
model~\cite{varga2026}. It never invokes equipment itself: each enabled step is handed to a WorkflowExecutor, which has
the provider orchestrated for itself and reports back. The WorkflowManager is generic, one instance per work-centre
local cloud, configured only by its workstation identifier.

\begin{figure}[ht!]
  \centering
  \resizebox{\textwidth}{!}{
\begin{tikzpicture}[node distance=4mm]
\node[ahsys,minimum height=5.25cm,minimum width=3.6cm] (sys) at (0,0) {WorkflowManager\\[-1pt]{\normalfont\scriptsize $\ll$system$\gg$}};
\node[ahsrv,anchor=west] (p0) at (3.4,2.25) {workflowRecipe};
\draw[thick,ArrowheadPurple] (1.8,2.25) -- (p0.west);
\fill[ArrowheadPurple] (1.8,2.25) circle (1.4pt);
\node[ahsrv,anchor=west] (p1) at (3.4,1.50) {workflowInstance};
\draw[thick,ArrowheadPurple] (1.8,1.50) -- (p1.west);
\fill[ArrowheadPurple] (1.8,1.50) circle (1.4pt);
\node[ahsrv,anchor=west] (p2) at (3.4,0.75) {workflowWorkItem};
\draw[thick,ArrowheadPurple] (1.8,0.75) -- (p2.west);
\fill[ArrowheadPurple] (1.8,0.75) circle (1.4pt);
\node[ahsrv,anchor=west] (p3) at (3.4,0.00) {workflowContext};
\draw[thick,ArrowheadPurple] (1.8,0.00) -- (p3.west);
\fill[ArrowheadPurple] (1.8,0.00) circle (1.4pt);
\node[ahsrv,anchor=west] (p4) at (3.4,-0.75) {workflowLease};
\draw[thick,ArrowheadPurple] (1.8,-0.75) -- (p4.west);
\fill[ArrowheadPurple] (1.8,-0.75) circle (1.4pt);
\node[ahsrv,anchor=west] (p5) at (3.4,-1.50) {workflowManagement};
\draw[thick,ArrowheadPurple] (1.8,-1.50) -- (p5.west);
\fill[ArrowheadPurple] (1.8,-1.50) circle (1.4pt);
\node[ahsrv,anchor=west] (p6) at (3.4,-2.25) {generalManagement};
\draw[thick,ArrowheadPurple] (1.8,-2.25) -- (p6.west);
\fill[ArrowheadPurple] (1.8,-2.25) circle (1.4pt);
\node[ahcsm,anchor=east] (c0) at (-3.4,2.25) {serviceDiscovery\\[-2pt]{\tiny\rmfamily ServiceRegistry}};
\draw[thick,ArrowheadTeal,-{Stealth[length=1.6mm]}] (c0.east) -- (-1.8,2.25);
\node[ahcsm,anchor=east] (c1) at (-3.4,1.50) {identity\\[-2pt]{\tiny\rmfamily Authentication}};
\draw[thick,ArrowheadTeal,-{Stealth[length=1.6mm]}] (c1.east) -- (-1.8,1.50);
\node[ahcsm,anchor=east] (c2) at (-3.4,0.75) {authorization\\[-2pt]{\tiny\rmfamily ConsumerAuthorization}};
\draw[thick,ArrowheadTeal,-{Stealth[length=1.6mm]}] (c2.east) -- (-1.8,0.75);
\node[ahcsm,anchor=east] (c3) at (-3.4,0.00) {workflowExecution\\[-2pt]{\tiny\rmfamily WorkflowExecutor}};
\draw[thick,ArrowheadTeal,-{Stealth[length=1.6mm]}] (c3.east) -- (-1.8,0.00);
\node[ahcsm,anchor=east] (c4) at (-3.4,-0.75) {workflowProduct\\[-2pt]{\tiny\rmfamily Intelligent Product}};
\draw[thick,ArrowheadTeal,-{Stealth[length=1.6mm]}] (c4.east) -- (-1.8,-0.75);
\node[ahcsm,anchor=east] (c5) at (-3.4,-1.50) {sysMonitor (AMQP)\\[-2pt]{\tiny\rmfamily SysMonitor broker}};
\draw[thick,ArrowheadTeal,-{Stealth[length=1.6mm]}] (c5.east) -- (-1.8,-1.50);
\node[font=\scriptsize\bfseries,ArrowheadPurple] at (5.0,3.15) {produced};
\node[font=\scriptsize\bfseries,ArrowheadTeal] at (-5.0,3.15) {consumed};
\end{tikzpicture}
  }
  \caption{SysML block view of the WorkflowManager: produced services (right) and consumed services with their providers (left).}
  \label{fig:block}
\end{figure}

The rest of this document is organized as follows.
In Section~\ref{sec:prior_art}, we reference major prior art capabilities of the system.
In Section~\ref{sec:use}, we describe the intended usage of the system.
In Section~\ref{sec:properties}, we describe fundamental properties provided by the system.
In Section~\ref{sec:delimitations}, we describe delimitations of capabilities of the system.
In Section~\ref{sec:services}, we describe the abstract service functions consumed or produced by the system.
In Section~\ref{sec:security}, we describe the security capabilities of the system.

\newpage
\subsection{Significant Prior Art}
\label{sec:prior_art}
Derhamy's Workflow Manager, a stateless book-keeper and blackboard that can be cloned per work
centre~\cite{derhamy2018,derhamy2018workflow}; Garcia Represa's workflow management system and the coexistence of
orchestration-style and choreography-style workflow management in one SoS~\cite{represa2022}; Varga's formal recipe
and session model with its transition rules and proven guarantees~\cite{varga2026}; and the Arrowhead Choreographer,
whose plans the WorkflowManager can compile recipes to.

\subsection{How This System Is Meant to Be Used}
\label{sec:use}
\begin{enumerate}
\item An engineer designs a workflow in the WorkflowConfigurator, which validates and registers it through
\texttt{workflowRecipe}.
\item A product, line controller or upper WorkflowManager starts a session through \texttt{workflowInstance}, either
naming the recipe or naming a product system that provides the recipe through \texttt{workflowProduct}.
\item The WorkflowManager enables steps, selects executors by their declared capabilities (push) or publishes work
items to be claimed (pull), and records bindings, results and failures from the executors' reports.
\item When every action has completed (Z-Close) it stores the results back at the product; the owner can read the full
trace at any time.
\end{enumerate}

\subsection{System functionalities and properties}
\label{sec:properties}
\subsubsection{Functional properties of the system}
\begin{itemize}
\item Recipe catalogue: validation (structure W1--W4, data flow, leases, guards, catalogue checks), immutable storage by
name and version, retirement, compilation to Choreographer plans.
\item Session engine: the session state $z = (M, M_A, \beta, \varepsilon, \xi, status)$ and the eight transition rules
A-Start, S-Enable, S-Bind, S-Done, A-Finish, Z-Close, S-FailBind and S-FailInvoke, with the Theorem~1 transition count.
\item Extensions, flagged per session: x-rebind (retry on ERROR or an unreachable executor), x-abort, x-timeout,
compensation (lease release on failure), x-guard/x-skip (guarded transitions with dead-path elimination).
\item Dispatch: push (executor selection by $\eta(s)$ and $\theta(s)$, least loaded, next executor after BIND\_FAILED),
pull (atomic claim, holder-only reports, claim timeout); \texttt{event} is accepted and handled as pull.
\item Session context ($\rho$, $\xi|dep(s)$) for executors; FIFO leases with renewal and expiry for shared resources.
\item Product integration (retrieve recipe at start, store results at Z-Close) and an audit worklog.
\end{itemize}

\subsubsection{Configuration of system properties}
The system is configured with Spring properties; the main ones are listed in the SysDD. The workflow-specific behaviour
is configured under the prefix \texttt{workflow.}; the security profile (TLS, authentication, authorization) is the
common Arrowhead~5 profile.

\subsubsection{Data stored by the system}
The WorkflowManager owns the database schema \texttt{ah\_workflow}: \texttt{recipe} (immutable canonical JSON per
version and its status), \texttt{workflow\_session} (the JSON snapshot of $z$ with its trace and an optimistic
\texttt{revision}), \texttt{work\_item} (one row per step attempt), \texttt{lease}, \texttt{worklog} and \texttt{logs}.
Storing $z$ as one snapshot keeps the stored state identical to the model object and lets sessions resume after a
restart.

\subsubsection{Non functional properties}
\begin{itemize}
\item \emph{Security:} see Section~\ref{sec:security}.
\item \emph{Safety and liveness:} for well-formed recipes without extensions every run takes at most $3|S| + 2|A|$
transitions and always progresses (Theorems~1--2 of~\cite{varga2026}); failure is absorbing. Exclusive tool use is
proved under assumption (A1) and holds for the lease variant only under an additional, validated configuration
assumption.
\item \emph{Latency:} one transition is a database update of the session snapshot; dispatch happens asynchronously on a
worker pool (\texttt{workflow.worker.threads}, default 8).
\item \emph{Energy and power saving:} no specific properties; the system is idle between requests except for the
claim/timeout watchdog.
\end{itemize}

\subsubsection{Stateful or stateless}
Stateful. Recipes, sessions (state $z$, trace), work items, leases and the worklog are persisted, so a
WorkflowManager restart resumes running sessions. Executors keep no durable state.

\subsection{Important Delimitations}
\label{sec:delimitations}
\begin{itemize}
\item No planning or scheduling (PlantDescriptionEngine or ERP), no quality policy authoring, no equipment services.
\item One local cloud: delegating actions to WorkflowManagers in other clouds through Gatekeeper/GateTunnel and the
disconnected working copy are designed but not implemented.
\item EventHandler notifications (\texttt{event} dispatch and workflow events) are not implemented.
\item Recovery beyond the listed extensions (compensation of arbitrary steps, \texttt{quantity} $>$ 1) is not provided;
the theorems do not cover sessions that use extensions.
\item The \texttt{workflowProduct} side is implemented as a consumer only; no product system exists yet.
\end{itemize}

\newpage
\section{Services}
\label{sec:services}
This section lists the produced and consumed services. Each produced service is specified by an SD and an HTTP/TLS/JSON
IDD in this documentation set.

\subsection{Produced services}

{\small
\begin{longtable}{|p{3.8cm}|p{6.2cm}|p{5.8cm}|}

\hline
\rowcolor{gray!33} Service & Operations & Documents \\ \hline
\endhead
\texttt{workflowRecipe} & validate, register, query, get, retire, compile & \texttt{WorkflowManager-v5.4.0-workflowRecipe-SD}, \texttt{-HTTP-IDD} \\ \hline
\texttt{workflowInstance} & start, abort, query, get-status, get-trace & \texttt{WorkflowManager-v5.4.0-workflowInstance-SD}, \texttt{-HTTP-IDD} \\ \hline
\texttt{workflowWorkItem} & claim, report & \texttt{WorkflowManager-v5.4.0-workflowWorkItem-SD}, \texttt{-HTTP-IDD} \\ \hline
\texttt{workflowContext} & get-parameters, get-results, put-result & \texttt{WorkflowManager-v5.4.0-workflowContext-SD}, \texttt{-HTTP-IDD} \\ \hline
\texttt{workflowLease} & reserve, poll, renew, release & \texttt{WorkflowManager-v5.4.0-workflowLease-SD}, \texttt{-HTTP-IDD} \\ \hline
\texttt{workflowManagement} & query-worklog, query-executors, query-work-items & \texttt{WorkflowManager-v5.4.0-workflowManagement-SD}, \texttt{-HTTP-IDD} \\ \hline
\texttt{generalManagement} & get-log, get-config & \texttt{Workflow-v5.4.0-generalManagement-SD}, \texttt{-HTTP-IDD} \\ \hline
\end{longtable}
}


\subsection{Consumed services}
All three workflow systems use the core systems the same way: at start-up they log in at the Authentication
system (\texttt{identity}) and register their services in the ServiceRegistry (\texttt{serviceDiscovery.register});
on every incoming request they verify the caller's identity token (\texttt{identity.identity-verify}) and, with
\texttt{service.authorization.enabled}, the caller's grant (\texttt{authorization.verify} of ConsumerAuthorization);
at shutdown they revoke their services and log out.

{\small
\begin{longtable}{|p{3.4cm}|p{3.2cm}|p{3.6cm}|p{5.6cm}|}

\hline
\rowcolor{gray!33} Service & Provider & Operations & Purpose \\ \hline
\endhead
\texttt{serviceDiscovery} & ServiceRegistry & register, revoke, lookup & Registration; discovery of executors (\texttt{workflowExecution} with its metadata) and the catalogue check C01. \\ \hline
\texttt{identity} & Authentication & login, logout, identity-verify & Own identity token; verification of callers. \\ \hline
\texttt{authorization} & ConsumerAuthorization & verify & Provider-side authorization of every service operation. \\ \hline
\texttt{workflowExecution} & WorkflowExecutor & start, abort & Push dispatch and abort of work items. \\ \hline
\texttt{workflowProduct} & Intelligent Product & retrieve-process-step, store-process-step-result & Product-driven sessions. Contract: \texttt{Workflow-v5.4.0-workflowProduct-SD}. \\ \hline
\texttt{sysMonitor (AMQP)} & SysMonitor broker & publish & Health and metrics when \texttt{sysmonitor.enabled}. \\ \hline
\end{longtable}
}


\newpage
\section{Security}
\label{sec:security}
The WorkflowManager runs in Arrowhead secure mode in every cloud deployment; an insecure mode exists only for development and
unit tests. It handles only Arrowhead-compliant X.509 certificates of its local cloud; others are refused during the
TLS handshake.

\subsection{Security model}
\begin{itemize}
\item \textbf{Secure mode only in deployment.} The deployment configuration (\texttt{deployment/*.properties}) enables
TLS with client authentication (\texttt{server.ssl.client-auth=need}, TLS 1.3 and 1.2), outsourced authentication with
identity tokens bound to the client certificate, interface policy \texttt{CERT\_AUTH} and the CRL check. The packaged
defaults without TLS exist for development and unit tests (\texttt{workflow.standalone}) and must not be used in a
cloud.
\item \textbf{Certificates.} The system has an Arrowhead system certificate
\texttt{<System>.<cloud>.<organisation>.arrowhead.eu} issued by the local cloud CA, kept in a PKCS\#12 key store; the
trust store holds the cloud CA. Non-Arrowhead certificates are refused at the TLS layer.
\item \textbf{Provider-side authorization.} Every operation of the registered services (other than management) is
checked against ConsumerAuthorization (\texttt{authorization.verify}) for the calling system, the service definition
and the operation (route tables \texttt{ServiceOperations} and \texttt{ConfiguratorOperations}). A missing grant gives
403.
\item \textbf{Management.} \texttt{generalManagement} and \texttt{workflowManagement} are protected by
\texttt{management.policy=sysop-only}.
\item \textbf{Data protection.} All service traffic is encrypted by TLS. The WorkflowManager's database connection uses
TLS with server certificate verification (\texttt{sslMode=verify-full}). Secrets (store, identity and database
passwords, the console password hash) come from container secrets (\texttt{configtree:/run/secrets/}), not from the
image.
\end{itemize}


\subsection{Object-level checks}
In addition to ConsumerAuthorization the WorkflowManager protects itself (first-person perspective): only the holder of
a work item may report on it or read and write its context; only the owner of a session or a sysop may abort it.

\newpage
\bibliographystyle{IEEEtran}
\bibliography{bibliography}

\newpage

\section{Revision History}
\subsection{Amendments}

\noindent\begin{tabularx}{\textwidth}{| p{1cm} | p{2.2cm} | p{1.6cm} | X | p{3.6cm} |} \hline
\rowcolor{gray!33} No. & Date & Version & Subject of Amendments & Author \\ \hline
1 & 2026-11-23 & \arrowversion & Initial edition: system description of the WorkflowManager, written from the Java reference implementation (ah-Workflow-java). & Documentation maintenance \\ \hline
\end{tabularx}

\subsection{Quality Assurance}

\noindent\begin{tabularx}{\textwidth}{| p{1cm} | p{2.2cm} | p{1.6cm} | X |} \hline
\rowcolor{gray!33} No. & Date & Version & Approved by \\ \hline
1 & & \arrowversion & Pending review \\ \hline
\end{tabularx}

\end{document}
